\clearpage\SourceLeaf{199}%
\pdfbookmark[1]{Gauss an Schumacher, 23. November 1847}{schumacher18471123}
\fancyhead[L]{\small GAUSS AN SCHUMACHER}
\begin{center}
{\small GAUSS AN SCHUMACHER.\par}
\medskip\hrule\medskip
{\small Astronomische Nachrichten herausgegeben von H.\ C.\ \textsc{Schumacher}.\par
Band 27.\ Nr.\ 625.\ S.\ 1.\par}
\medskip\hrule
\end{center}\medskip
\begin{flushright}Göttingen, 23. November 1847.\end{flushright}
Die in Nr.\ 615 der A.\ N.\ abgedruckten zweiten Bahnelemente der Iris hatte Herr Prof.\ \textsc{Goldschmidt} gleich nach dem 19.\ September berechnet. Sie haben auch noch in den folgenden Monaten gut mit den Beobachtungen übereingestimmt (am 17.\ und 18.\ November betrug die Abweichung von Herrn \textsc{Rümcker}s Meridianbeobachtung $16''$ in Rektaszension und $2''$ in Deklination). Man darf sie daher bereits als sehr gute Näherung betrachten. Ich habe deshalb Herrn Prof.\ \textsc{Goldschmidt} veranlasst, damit den Zodiakus der Iris zu berechnen. Das Ergebnis dieser Arbeit lasse ich hier folgen. Es kann dazu dienen, leichter zu prüfen, ob frühere Beobachtungen denselben Himmelskörper betrafen.

Bei dieser Gelegenheit nehme ich den Aufsatz wieder zur Hand, in dem ich vor 44 Jahren die allgemeine Methode zur Bestimmung der Grenzen eines solchen Zodiakus angegeben habe. Ich sehe, dass ich darin auch schon die Ausnahmefälle angedeutet habe, in denen der Bereich auf der Himmelskugel, in dem ein nach \textsc{Kepler}s Gesetzen um die Sonne kreisender Himmelskörper geozentrisch erscheinen kann, entweder gar keine oder nur \textit{eine} Grenze hat. Dabei ergibt die Methode im ersten Fall eine in sich zurücklaufende Linie und im anderen zwei solche Linien. Dort ist auch die Frage aufgeworfen, was in solchen Ausnahmefällen diese durch Rechnung gefundenen Linien eigentlich
\SourceLeaf{200}%
bedeuten. Ich habe mich damals auf diese Andeutungen beschränkt, weil eine weitere Ausführung dort ein Horsd’œuvre gewesen wäre und ich anderen gern das Vergnügen lassen wollte, sich mit einer meiner Meinung nach nicht uninteressanten mathematischen Aufgabe zu beschäftigen. Da ich aber von niemandem erfahren habe, der in der langen Zwischenzeit die mir inzwischen ganz entfallene Frage aufgegriffen hätte, nutze ich diese Gelegenheit, um hier wenigstens den Hauptnerv dessen mitzuteilen, was zur Beantwortung nötig ist.

Eine geozentrische Position des betreffenden Planeten (oder Kometen) ergibt sich, indem man einen Punkt seiner Bahn mit einem Punkt der Erdbahn kombiniert. Doch kann sich dieselbe geozentrische Position auch aus zwei, drei oder vier verschiedenen Kombinationen ergeben. Um sich das anschaulich zu machen, lege man durch die Sonne eine Ebene, auf der die Gerade senkrecht steht, die eine vorgegebene geozentrische Position auf der Himmelskugel darstellt. Auf diese Ebene projiziere man sowohl die Erdbahn als auch die Planetenbahn orthografisch. Beide Projektionen sind Ellipsen oder allgemeiner Kegelschnitte, die sich entweder gar nicht, zweimal oder viermal schneiden. Eine Berührung ist dabei als Verschmelzen zweier Schnittpunkte anzusehen. Jeder Schnittpunkt stellt eine geozentrische Position dar, die entweder mit der vorgegebenen Position übereinstimmt oder ihr auf der Himmelskugel gegenüberliegt. Das hängt davon ab, ob von den beiden Punkten auf der Planetenbahn und der Erdbahn, deren Projektionen zusammenfallen, der erste oder der zweite höher liegt. Als obere Seite der Projektionsebene gilt die Seite, auf der die vorgegebene geozentrische Position liegt.

Daraus ist klar, dass, wenn man irgendeinen Punkt der Himmelskugel. versuchsweise als geozentrische Position vorgibt, dieser es entweder auf gar keine, auf \textit{eine} oder auf zwei, drei oder vier Arten sein kann. Für einen gegebenen Planeten teilt sich so die ganze Oberfläche der Himmelskugel in verschiedene Bereiche. Die nach der von mir angegebenen Methode bestimmten Linien sind, allgemein gesprochen, lediglich Trennlinien zwischen zwei Bereichen der Kugeloberfläche, bei denen die Punkte auf der einen Seite auf zwei Arten mehr geozentrische Positionen sein können als auf der anderen. Die Punkte auf den Trennlinien bilden den Übergang. Das heißt, sie können auf \textit{eine} Art mehr als die Punkte auf der einen Seite und auf \textit{eine} Art weniger als die
\SourceLeaf{201}%
Punkte auf der anderen Seite geozentrische Positionen sein. Übrigens lässt sich auch das Kriterium angeben, nach dem man a priori entscheidet, auf welcher Seite der Trennlinie es zwei Lösungen mehr gibt als auf der anderen. Dabei möchte ich mich aber jetzt nicht aufhalten.

Alle bisher bekannten Planeten haben solche Bahnen, dass die durch die Theorie gefundenen Trennlinien immer wirkliche Grenzen sind. Es sind nämlich zwei Trennlinien, die die Himmelskugel in drei Bereiche teilen. Zwei davon sind getrennt liegende Flächen, in die gar keine geozentrischen Positionen fallen. Die dritte liegt gürtelförmig dazwischen und enthält alle geozentrischen Positionen. Jeder Punkt innerhalb des Gürtels kann auf zwei Arten geozentrische Position sein, jeder Punkt auf der Grenze auf \textit{eine} Art.

Man kann sich aber Bahnen vorstellen (mehrere Kometenbahnen fallen vielleicht darunter, doch habe ich das bisher nicht untersucht), bei denen zwei Grenzen die Himmelskugel ebenfalls in drei Bereiche teilen. Der eine Bereich enthält gar keine geozentrischen Positionen, der andere die geozentrischen Positionen auf zwei Arten und der dritte die Punkte, die auf vier verschiedene Arten geozentrische Positionen sein können.

Weitere Vielfalt ergibt sich, wenn \textit{eine} Trennlinie sich einmal oder zweimal selbst schneidet und so eine einfache oder doppelte Schleife bildet. Im ersten Fall trennt sie zwei, im zweiten drei Bereiche vom gürtelförmigen Teil ab, in denen es entsprechend 4 und 0 oder 4, 0, 4 Lösungsarten gibt.

\Needspace{8\baselineskip}
Wieder anders verhält es sich mit einer Bahn, die wie ein Kettenglied in die Erdbahn eingreift. In einem solchen Fall gibt es nur \textit{eine} zusammenhängende Grenzlinie und daher zwei getrennte Flächenbereiche. Die Punkte des einen können auf \textit{eine} Art, die des anderen auf drei Arten geozentrische Positionen sein.
\par\medskip\centerline{\rule{.25\textwidth}{.3pt}}
